%%%%%%%%%%%%%%%%%%%%%%%%%%%%%%%
%%% Copier-Coller MACINTOSH %%%
%%%%%%%%%%%%%%%%%%%%%%%%%%%%%%%
\index{Copier}
\index{Coller}
\begin{figure}[htbp]
\begin{center}
\begin{tabular}{@{}|p{13.2cm}|@{}} \hline
\multicolumn{1}{|c|}{\raisebox{-2ex}{\rule{0ex}{5ex}}\em Sélectionner et recopier un texte sur Macintosh}
\\  \hline \hline
{% debut de la mini-page
\begin{minipage}[t]{13.2cm}
\vspace{2ex}
\begin{enumerate}
\item Placer le curseur d'édition clignotant au début du texte à 
recopier en le
déplaçant à l'aide des flèches de défilement
(\hspace{0.3ex}\begin{picture}(12,12)(0, 2)
\put(0,0){\framebox(12,12)[l]{$\:\scriptstyle\uparrow$}}
\end{picture}\hspace{0.3ex},
\begin{picture}(12,12)(0, 2)
\put(0,0){\framebox(12,12)[l]{$\:\scriptstyle\downarrow$}}
\end{picture}\hspace{0.3ex},
\begin{picture}(12,12)(0, 2)
\put(0,0){\framebox(12,12)[bl]{$\:\scriptstyle\leftarrow$}}
\end{picture}\hspace{0.3ex},
\hbox{\begin{picture}(12,12)(0, 2)
\put(0,0){\framebox(12,12)[bl]{$\:\scriptstyle\rightarrow$}}
\end{picture}\hspace{0.3ex})} ou bien en cliquant avec la souris
à l'endroit désiré.
\item Maintenir la touche majuscules 
\framebox[11.mm][l]{\rule{0mm}{2.5mm}
\raisebox{-0.75mm}[0mm][0mm]{\includevspacePostScript[3mm]{Shift}{0.95}{0.95}}} 
ou  \hspace{0.3ex}\framebox[11.mm][l]{~~~\tiny\texttt{shift}}\hspace{0.3ex}
enfoncée, sans la relâcher.
\item Placer le curseur d'édition clignotant à la fin du texte 
à recopier en le
déplaçant à l'aide des flèches de défilement
(\hspace{0.3ex}\begin{picture}(12,12)(0, 2)
\put(0,0){\framebox(12,12)[l]{$\:\scriptstyle\uparrow$}}
\end{picture}\hspace{0.3ex},
\begin{picture}(12,12)(0, 2)
\put(0,0){\framebox(12,12)[l]{$\:\scriptstyle\downarrow$}}
\end{picture}\hspace{0.3ex},
\begin{picture}(12,12)(0, 2)
\put(0,0){\framebox(12,12)[bl]{$\:\scriptstyle\leftarrow$}}
\end{picture}\hspace{0.3ex},
\hbox{\begin{picture}(12,12)(0, 2)
\put(0,0){\framebox(12,12)[bl]{$\:\scriptstyle\rightarrow$}}
\end{picture}\hspace{0.3ex})} ou bien en cliquant avec la souris
à l'endroit désiré, de sorte que le texte à recopier
passe en {\em vidéo-inverse} c'est-à-dire en blanc sur fond noir
(si par erreur vous tapez du texte au clavier, il remplace le texte
sélectionné en vidéo-inverse$\:$; la touche d'effacement 
\framebox[11.mm][l]{$\leftarrow$~~~}
supprime le texte sélectionné en vidéo-inverse$\:$; dans les deux cas,
choisir l'option {\idt Annuler}
du menu {\idt Edition} pour rétablir l'original).
%
\item Relâcher la touche majuscules
\framebox[11.mm][l]{\rule{0mm}{2.5mm}
\raisebox{-0.75mm}[0mm][0mm]{\includevspacePostScript[3mm]{Shift}{0.95}{0.95}}}$\,$ ou  \hspace{0.3ex}\framebox[11.mm][l]{~~~\tiny\texttt{shift}}\hspace{0.3ex}.
\item Sélectionner l'option {\idt Copier} du menu {\idt Edition}
(ou de manière équivalente maintenir
la touche commande
\framebox[9.mm][l]{\rule{0mm}{2.5mm}
\raisebox{-0.5mm}[0mm][0mm]{\includevspacePostScript[3mm]{Cmd}{0.95}{0.95}}}
 enfoncée et taper sur la touche 
\begin{picture}(12,12)(0, 2) %touche d'echappement
\put(0,0){\framebox(12,12)[lb]{}}
\put(1.5,2){\makebox(0,0)[lb]{\tiny C}}
\end{picture}\hspace{0.3ex}).
%
\item Placer le curseur d'édition clignotant à l'endroit où le texte 
doit être recopié en le
déplaçant à l'aide des flèches de défilement
(\hspace{0.3ex}\begin{picture}(12,12)(0, 2)
\put(0,0){\framebox(12,12)[l]{$\:\scriptstyle\uparrow$}}
\end{picture}\hspace{0.3ex},
\begin{picture}(12,12)(0, 2)
\put(0,0){\framebox(12,12)[l]{$\:\scriptstyle\downarrow$}}
\end{picture}\hspace{0.3ex},
\hbox{\begin{picture}(12,12)(0, 2)
\put(0,0){\framebox(12,12)[bl]{$\:\scriptstyle\leftarrow$}}
\end{picture}\hspace{0.3ex},}
\begin{picture}(12,12)(0, 2)
\put(0,0){\framebox(12,12)[bl]{$\:\scriptstyle\rightarrow$}}
\end{picture}\hspace{0.3ex}) ou bien en cliquant à l'endroit
désiré avec la souris.
\item Sélectionner l'option {\idt Coller} du menu {\idt Edition}
(ou de manière équivalente maintenir
la touche commande 
\framebox[9.mm][l]{\rule{0mm}{2.5mm}
\raisebox{-0.5mm}[0mm][0mm]{\includevspacePostScript[3mm]{Cmd}{0.95}{0.95}}}
enfoncée et taper sur la touche 
\begin{picture}(12,12)(0, 2) %touche d'echappement
\put(0,0){\framebox(12,12)[lb]{}}
\put(1.5,2){\makebox(0,0)[lb]{\tiny V}}
\end{picture}\hspace{0.3ex}).
\end{enumerate}
L'opération de recopie ({\idt Coller}) en un endroit quelconque du
programme peut être répétée autant de
fois que nécessaire.
\vspace{2ex}
\end{minipage} 
} % fin de la mini-page
\\ \hline
\end{tabular}
\end{center}
\caption{Copier-coller sur le Macintosh}\label{fig:mac-copier-coller}
\end{figure}